% ============================================================
% PART II: Generalisations — summary slides 14-18
% ============================================================

\section{Part II: Generalisations}

% ----------------------------------------------------------
% SLIDE 14: Part II opening
% SAY: "Part I assumed unpolarised illumination and used T_0 as
%       a constant. What happens when we relax these two
%       assumptions? Arbitrary polarisation is handled by the
%       absorption Stokes vector. Sub-6 GHz is handled by the
%       exact flux-averaged transmission T-bar."
% ----------------------------------------------------------
\begin{frame}{Part II: relaxing two assumptions}
\hypertarget{sec:part2}{}

  \begin{columns}[T]
    \begin{column}{0.48\textwidth}
      Part I assumed:
      \begin{enumerate}
        \item Unpolarised or circular illumination ($q = 0$)
        \item mmWave frequencies ($\Tavg \approx T_0$)
      \end{enumerate}

      \vspace{8pt}
      Part II retains both exactly:
      \begin{enumerate}
        \item Arbitrary polarisation\\
              $\rightarrow$ absorption Stokes vector
        \item Any frequency above ${\sim}\SI{100}{\mega\hertz}$\\
              $\rightarrow$ exact flux-averaged $\Tbar$
      \end{enumerate}
    \end{column}

    \begin{column}{0.48\textwidth}
      \vspace{12pt}
      {\small
      Key results:
      \begin{itemize}
        \item Worst-case polarisation correction on human body: $16\,\%$
        \item Universal bound (cylinder): $27.9\,\%$
        \item With $\ge 20$ multipath: $< 2.5\,\%$
        \item $\Tbar$ makes Cauchy formula exact at any frequency
      \end{itemize}
      }
    \end{column}
  \end{columns}
\end{frame}

% ----------------------------------------------------------
% SLIDE 15: Absorption Stokes vector
% SAY: "Total absorbed power is an inner product of the incident
%       Stokes vector with a body response vector. The body
%       response is precomputed from the mesh. Fourth component
%       is zero: circular handedness does not matter. Worst case
%       on the body is 16%, bounded by the cylinder at 27.9%."
% ----------------------------------------------------------
\begin{frame}{Absorption Stokes vector}
\hypertarget{sum-p2-stokes}{}

  Retain polarisation exactly. Integrate over body surface:

  \begin{resultbox}[Stokes vector representation]
    \vspace{-6pt}
    \[
      P_{\mathrm{abs}} = \mm(\khat) \cdot \ssv_{\mathrm{inc}}
    \]
  \end{resultbox}

  \vspace{2pt}
  \begin{columns}[T]
    \begin{column}{0.52\textwidth}
      $\ssv_{\mathrm{inc}} = (S_0, S_1, S_2, S_3)^\top$: incident Stokes vector

      \vspace{4pt}
      $\mm(\khat) = \bigl(\tilde{A}_\perp,\; -\tfrac{1}{2}B_c,\; -\tfrac{1}{2}B_s,\; 0\bigr)^\top$

      \vspace{4pt}
      \begin{itemize}
        \item Both $\tilde{A}_\perp$ and $\tilde{B}$ precomputed from mesh
        \item Fourth component $= 0$: circular handedness irrelevant
      \end{itemize}
      {\footnotesize \gotodetail{app-p2-stokes}{Stokes derivation}}
    \end{column}

    \begin{column}{0.44\textwidth}
      Polarisation directivity $D_B$:
      \begin{itemize}
        \item Bounds fractional correction
        \item Cylinder (worst case): $27.9\,\%$
        \item Thelonious (max over $S^2$): $16.0\,\%$
      \end{itemize}

      \vspace{4pt}
      {\footnotesize Unpolarised geometric law is conservative
       for vertical polarisation: horizontal limbs are purely TE,
       $T_s \le T_0$.\\
       \gotodetail{app-p2-sensitivity}{$D_B$ bounds}\;
       \gotodetail{app-p2-canonical}{Canonical shapes}}
    \end{column}
  \end{columns}
\end{frame}

% ----------------------------------------------------------
% SLIDE 16: Multipath convergence
% SAY: "The standard deviation of the polarisation error scales
%       as D_B over root 2N. With 20 paths, that is below 2.5%,
%       comparable to the pseudo-Brewster approximation itself.
%       In any realistic indoor environment, polarisation is
%       a non-issue for total absorbed power."
% ----------------------------------------------------------
\begin{frame}{Multipath convergence}
\hypertarget{sum-p2-multipath}{}

  \begin{columns}[T]
    \begin{column}{0.48\textwidth}
      With $N$ independent paths, polarisation errors cancel:

      \[
        \frac{\sigma}{P_{\mathrm{abs}}} = \frac{D_B}{\sqrt{2N}}
      \]

      \vspace{6pt}
      \begin{tabular}{rc}
        \toprule
        $N$ paths & Pol.\ error (std) \\
        \midrule
        1  & up to $16\,\%$ \\
        10 & $\le 3.6\,\%$ \\
        20 & $\le 2.5\,\%$ \\
        50 & $\le 1.6\,\%$ \\
        \bottomrule
      \end{tabular}

      \vspace{6pt}
      {\footnotesize In realistic indoor mmWave channels ($N \ge 20$),
       polarisation error is below the pseudo-Brewster error.\\
       \gotodetail{app-p2-multipath}{Convergence details}}
    \end{column}

    \begin{column}{0.48\textwidth}
      \begin{center}
        \only<1>{%
          \includegraphics[width=\textwidth]{multipath_panel_a.png}\\[2pt]
          {\footnotesize (a) Convergence of polarisation error with $N$.\\
           $D_B/\sqrt{2N}$ bound (dashed).}
        }%
        \only<2->{%
          \includegraphics[width=\textwidth]{multipath_panel_b.png}\\[2pt]
          {\footnotesize (b) Histogram of $D_B$ over directions.\\
           Max $16\,\%$, mean $8\,\%$, cylinder bound $27.9\,\%$.}
        }%
      \end{center}
    \end{column}
  \end{columns}
\end{frame}

% ----------------------------------------------------------
% SLIDE 17: Sub-6 GHz extension
% SAY: "T-bar is the flux-averaged transmission. It replaces T_0
%       in the Cauchy formula and makes it exact from 100 MHz to
%       100 GHz. Below 40 GHz, the T_0 formula underestimates
%       (conservative). Below 6 GHz the local surface map loses
%       physical meaning, but total power via T-bar remains valid."
% ----------------------------------------------------------
\begin{frame}{Sub-\SI{6}{\giga\hertz} extension}
\hypertarget{sum-p2-sub6}{}

  \begin{columns}[T]
    \begin{column}{0.48\textwidth}
      {\small $\Tavg(\theta) \approx T_0$ weakens below \SI{10}{\giga\hertz}.}
      \vspace{0pt}
      {\small Replace $T_0$ with the exact flux-averaged transmission:}
      \vspace{-4pt}
      \[
        \Tbar(f) = 2\int_0^1 \Tavg(\mu, f)\;\mu\;\diff\mu
      \]
      \vspace{-6pt}
      \begin{resultbox}[Exact at any frequency]
        \vspace{-6pt}
        \[
          \langle P_{\mathrm{abs}} \rangle = \Sinc \, \Tbar \, \frac{\Aab}{4}
        \]
        \vspace{-6pt}
      \end{resultbox}
      \vspace{0pt}
      {\footnotesize
        Below \SI{40}{\giga\hertz}: $T_0$ underestimates (conservative).\\
        Below \SI{6}{\giga\hertz}: local map loses meaning;
        total power via $\Tbar$ remains valid.
      }
    \end{column}
    \begin{column}{0.48\textwidth}
      \begin{center}
        \only<1>{%
          \includegraphics[width=\textwidth]{R_of_f_panel_a.png}\\[2pt]
          {\footnotesize (a) Pseudo-Brewster ratio $R(f) = T_0/\Tbar$.\\
           Conservative ($R < 1$) below \SI{40}{\giga\hertz}.}
        }%
        \only<2>{%
          \includegraphics[width=\textwidth]{R_of_f_panel_b.png}\\[2pt]
          {\footnotesize (b) $T_0$ vs $\Tbar$ across frequency.\\
           Gap widens below \SI{6}{\giga\hertz}.}
        }%
        \only<3->{%
          \includegraphics[width=\textwidth]{R_of_f_panel_c.png}\\[2pt]
          {\footnotesize (c) $\Tavg(\theta)\cos\theta$ at representative frequencies.\\
           \gotodetail{app-p2-tbar}{$\Tbar$ derivation}}
        }%
      \end{center}
    \end{column}
  \end{columns}
\end{frame}

% ----------------------------------------------------------
% SLIDE 18: Part II detail index
% ----------------------------------------------------------
\begin{frame}{Part II: detail index}
  \label{p2-index}

  \begin{columns}[T]
    \begin{column}{0.48\textwidth}
      \gotodetail{app-p2-stokes}{Stokes vector derivation}\\[4pt]
      \gotodetail{app-p2-sensitivity}{Polarisation sensitivity $D_B$}\\[4pt]
      \gotodetail{app-p2-canonical}{Canonical shapes (sphere, cylinder)}\\[4pt]
      \gotodetail{app-p2-body-pol}{Polarisation on human body}\\[4pt]
    \end{column}

    \begin{column}{0.48\textwidth}
      \gotodetail{app-p2-multipath}{Multipath convergence details}\\[4pt]
      \gotodetail{app-p2-sh}{SH compression of polarisation}\\[4pt]
      \gotodetail{app-p2-tbar}{Exact averaged transmission $\Tbar$}\\[4pt]
      \gotodetail{app-p2-validity}{Validity regimes}\\[4pt]
    \end{column}
  \end{columns}
\end{frame}
